\begin{claim}{F.31}{Proved}{Finite affine marginal image equality on the full product-one character locus}
Use the actual bases $Y_i^-=X_i$ and
$Y_i^+=\mathcal B_+(X_i)$ under the based
isomorphism of F.23.13.
For every finite field $k$ and continuous
$\rho:V_\pm\to\operatorname{Aff}_1(k)$ whose
scalar character $\chi$ belongs to the exact locus
\[
                         \prod_{i=1}^7\chi(Y_i^\pm)=1,
                                                               \tag{F.31.1}
\]
the actual marginal images satisfy
\[
                         \rho(\mathcal K_\pm)=\rho(R_\pm).
                                                               \tag{F.31.2}
\]
All characteristics, all scalar-field extension
degrees and all larger original coefficient
fields are included.

Put $C=\chi(R_\pm)$,
$k_1=\mathbb F_p(C)$ and
$k_0=\mathbb F_p(\chi(V_\pm))$.
On the part $C\ne\{1\}$, set
$q_1=|k_1|$, $m=[k_0:k_1]$ and
$\chi_j=\chi^{q_1^j}$ for $0\leq j<m$.
The summed continuous $H^1$ restriction
for these $m$ characters is injective.
The corresponding map on cocycles has
exactly the principal tuples whose
parameters sum to zero as its kernel,
and has rank $6m+1$ on $7m$ coordinates.
\end{claim}

\begin{proof}
\emph{Actual scalar fields and completed scalar cases.}
Work first with
$K=\mathcal K_-\leq R=R_-\leq V=V_-$.
The negative character has the unique form
\[
 \begin{gathered}
 \chi=(r,s,t,u,v,w,c^{-1}),\qquad c=rstuvw,\\
 C=\chi(K)=\chi(R)=\langle u,v,w,c^{-1}\rangle,\\
 k_1=\mathbb F_p(u,v,w,c),\qquad
 k_0=k_1(r,t),\qquad s=\frac{c}{rtuvw}.
 \end{gathered}                                                \tag{F.31.3}
\]
All six displayed parameters are nonzero.
Here the equality of the scalar images is
the image of
$\overline{\langle\!\langle K\rangle\!\rangle}_R=R$
under the actual abelian scalar homomorphism;
F.21 and the marginal surjection of F.22
give precisely this normal-generation equality.
The actual deleted generators $X_4,\ldots,X_7$
give the displayed generating set for $C$.
The two descriptions of $k_0$ follow from
the last identity in (F.31.3), in both directions.

For $C=\{1\}$ the image of $R$ is an
abelian translation group, and the same
normal-generation map proves (F.31.2).
For $C\ne\{1\}$ and $k_0=k_1$, F.25
gives the image equality, while F.24
gives the asserted individual restriction
injection and cocycle kernel.
For the rest of the proof take
$C\ne\{1\}$ and $m>1$.

\emph{Coefficient maps and actual source actions.}
Let $P_0=X_1\cdots X_7$.
The coefficient homomorphism on the original
seven-variable derivations is exactly
\[
 \begin{aligned}
 \mathbb Z[t_1^{\pm1},\ldots,t_7^{\pm1}]&\longrightarrow
 A=\mathbb Z[r^{\pm1},s^{\pm1},t^{\pm1},u^{\pm1},v^{\pm1},w^{\pm1}],\\
 (t_1,t_2,t_3,t_4,t_5,t_6,t_7)&\longmapsto
          (r,s,t,u,v,w,(rstuvw)^{-1}).
 \end{aligned}                                                \tag{F.31.4}
\]
An exponent vector $(a_i)$ therefore maps to
$r^{a_1-a_7}s^{a_2-a_7}t^{a_3-a_7}
 u^{a_4-a_7}v^{a_5-a_7}w^{a_6-a_7}$.
The second chart applies the additional
coefficient homomorphism $r\mapsto1$.
Every computation below is an integral
Laurent identity on these specified rings.

These derivations are evaluated on the actual
discrete free words, with
$d(xy)=d(x)+\chi(x)d(y)$ and
$d(x^{-1})=-\chi(x)^{-1}d(x)$.
Every seven-tuple of values in $k_0$ gives
a homomorphism from the discrete free group
to its finite affine group, hence extends
uniquely and continuously to $V=\widehat F_7$.
The continuous cocycle space is thus exactly
$k_0^7$.
Its principal line is generated by
$\delta=(r,s,t,u,v,w,c^{-1})^{\mathsf T}-\mathbf1$.

Retain the eleven original F.23 automorphisms
$B_l(x)=u_l\alpha_l(x)u_l^{-1}$,
where throughout this proof $0\leq l\leq10$
is the original zero-based source index.
Thus $l=0$ is \texttt{old0}; indices $1,\ldots,5$
are \texttt{old\_1\_0},\ldots,\texttt{old\_1\_4},
and indices $6,\ldots,10$ are
\texttt{old\_-1\_0},\ldots,\texttt{old\_-1\_4}.
The source position in a one-based list is
always $l+1$.
Both $B_l$ and its inverse preserve $K$.
The pure Artin action fixes every abelian
character and the actual boundary word.
Write $A_l$ for precomposition on the seven
cocycle coordinates. Then
\[
 \begin{gathered}
 A_l\delta=\delta,\qquad
 D=(1,r,rs,rst,rstu,rstuv,c),\qquad D\delta=0,\\
 DA_l=\mu_lD,\qquad \mu_l=\chi(u_l),\\
 (\mu_0,\ldots,\mu_{10})=
 (1,1,rt,1,1,1,u^{-2},t^{-1},v^{-2},w^{-2},c^2).
 \end{gathered}                                                \tag{F.31.5}
\]
For the boundary identity, the exact cocycle
formula gives
$d(B_l(P_0))=\chi(u_l)d(P_0)$,
because $\chi(P_0)=1$ and $\alpha_l(P_0)=P_0$.
The displayed character values are evaluated
from the actual decorations in the eleven
labelled Artin words.

The word recurrence used for the finite
certificate retains both each meridian
image $x_i$ and its decoration $a_i$.
At a positive crossing of positions $j,j+1$ it sends
$(x_j,x_{j+1};a_j,a_{j+1})$ to
$(x_jx_{j+1}x_j^{-1},x_j;x_ja_{j+1},a_j)$;
at a negative crossing it sends them to
$(x_{j+1},x_{j+1}^{-1}x_jx_{j+1};
 a_{j+1},x_{j+1}^{-1}a_j)$.
The eighth meridian is evaluated at the
identity. At the end $u_l=a_8^{-1}$,
and its conjugation gives the seven rows
of $A_l$. The certificate checks the
label permutation, the erased meridian,
all seven preserved character values and
the full boundary product after each word.
These are the original source maps used
in every matrix below.

\emph{The retained principal chart for $r\ne1$.}
The principal quotient and its section are
\[
 \begin{aligned}
 P_1(z)&=(z_i-\delta_i(r-1)^{-1}z_1)_{i=2}^7,
 &S_1(x)&=(0,x_1,\ldots,x_6)^{\mathsf T},\\
 P_1S_1&=I_6,
 &S_1P_1(z)&=z-\delta(r-1)^{-1}z_1.
 \end{aligned}                                                \tag{F.31.6}
\]
They identify
$Q_\chi=H^1_{\rm cont}(V,k_{0,\chi})$
with the coordinates of the original
$X_2,\ldots,X_7$ in that order.
The inverse formulas prove the identification
and its exact kernel, namely $k_0\delta$.
Put $b=DS_1=(r,rs,rst,rstu,rstuv,c)$.
The boundary quotient $\pi(z)=bz$ is
well-defined since $D\delta=0$.
Its kernel is the invariant subspace $W_\chi$
of dimension five, using the explicit maps
\[
 \begin{gathered}
 \iota(x_1,\ldots,x_5)=
 (-sx_1-stx_2-stux_3-stuvx_4-stuvwx_5,x_1,\ldots,x_5),\\
 \sigma(a)=r^{-1}ae_1,\qquad
 \pi\sigma=I_1,\qquad
 z=\iota(z_2,\ldots,z_6)+\sigma\pi(z),\\
 J=(e_2-se_1,e_3-ste_1,e_4-stue_1,
              e_5-stuve_1,e_6-stuvwe_1,r^{-1}e_1),\\
 J^{-1}(z)=(z_2,z_3,z_4,z_5,z_6,bz).
 \end{gathered}                                                \tag{F.31.7}
\]
The displayed decomposition proves both
inverse basis identities; it does not
require $\sigma(k_0)$ to be invariant.
In this basis the actual source block is
\[
 T_l=J^{-1}P_1A_lS_1J=
       \begin{pmatrix}G_l&*\\0&\mu_l\end{pmatrix},
                    \qquad G_0=rI_5.
                                                               \tag{F.31.8}
\]
The scalar identity for $G_0$ is also the
literal-word identity (F.30.9).
The present certificate checks every entry
of all eleven blocks directly from the
original words. Its exact quotient check is
\[
 (r-1)(A_lS_1J-S_1JT_l)
       =\delta\,\operatorname{row}_1(A_lS_1J).
                                                               \tag{F.31.9}
\]
The section has zero first row, so applying
$P_1$ to this identity and cancelling
the proved nonzero $r-1$ gives exactly
(F.31.8), with no coordinate substitution.

\emph{Eight exact moving-root rows on $W_\chi$.}
The following table specifies a list of base
roots $a_{l,h}$ and one moving root $b_l$.
Let $Q_l(Z)=\prod_h(Z-a_{l,h})$.
Repeated base roots are retained as repeated
factors of this polynomial.
\[
 \begin{array}{c|c|c|c}
 l&\text{source position}&(a_{l,h})_h&b_l\\\hline
 6&7&(1,(vw)^{-1})&c/t\\
 8&9&(1,c/u)&(tw)^{-1}\\
 9&10&(1,c/v)&(tu)^{-1}\\
 10&11&(1,(uw)^{-1})&(tv)^{-1}\\
 1&2&(1,(vw)^{-1},(r^2u^2)^{-1})&c/t\\
 3&4&(1,c/u,(r^2v^2)^{-1})&(tw)^{-1}\\
 4&5&(1,c/v,(r^2w^2)^{-1})&(tu)^{-1}\\
 5&6&(1,(uw)^{-1},(c/r)^2)&(tv)^{-1}
 \end{array}.
                                                               \tag{F.31.10}
\]
For the displayed rows below the exact
integral identities are
\[
 \begin{gathered}
 Q_l(G_l)(G_l-b_lI_5)=0,\qquad q_lG_l=b_lq_l,\\
 q_lQ_l(G_l)=\tau_lq_l,
       \qquad \tau_l=Q_l(b_l)=\prod_h(b_l-a_{l,h}).
 \end{gathered}                                                \tag{F.31.11}
\]
The last equality follows by polynomial
evaluation from the second, and is also
checked directly. No conclusion about
diagonalizability is needed.

Here are all the rows used in (F.31.11).
To keep their coordinates explicit, set
$A=rst$, $B=t-1$, $H=rt-t+1$ and
$L=ru+(r-1)v-r$. These are the displayed
Laurent expressions, with $c=Auvw$.
Each displayed column is the transpose
of the named row:
\[
 \begin{gathered}
 q_6^{\mathsf T}=
 \begin{pmatrix}
 -A(c-1)H\\-B(c-1)(AH+1)\\
 -AuB(c-1)(H-u)\\-AuB(c-1)\bigl(v(H-u)+u\bigr)\\
 -cB\bigl((c-1)(H-u)+cu/v\bigr)
 \end{pmatrix},\qquad
 q_8^{\mathsf T}=
 \begin{pmatrix}w-1\\B(w-1)\\uB(w-1)\\-uvB\\0\end{pmatrix}.
 \end{gathered}                                                \tag{F.31.12}
\]
\[
 \begin{gathered}
 q_9^{\mathsf T}=
 \begin{pmatrix}
 A(u-1)\\B\bigl(A(u-1)-1\bigr)\\AuB(u-1)\\AuvB(u-1)\\cB(u-1)
 \end{pmatrix},\qquad
 q_{10}^{\mathsf T}=
 \begin{pmatrix}
 v-1\\B(v-1)\\-uB\bigl((r-1)v+1\bigr)\\
 -uv(r-1)B(v-1)\\-uvw(r-1)B(v-1)
 \end{pmatrix}.
 \end{gathered}                                                \tag{F.31.13}
\]
\[
 \begin{gathered}
 q_1^{\mathsf T}=
 \begin{pmatrix}c-1\\B(c-1)\\uB(c-1)\\uvB(c-1)\\cuvwB\end{pmatrix},\qquad
 q_3^{\mathsf T}=
 \begin{pmatrix}
 stv(w-1)\\B(w-1)(stv+1)\\-stuv(r-1)B(w-1)\\
 -stuvB\bigl((w-1)L+ru\bigr)\\-stuvwB(w-1)L
 \end{pmatrix}.
 \end{gathered}                                                \tag{F.31.14}
\]
\[
 \begin{gathered}
 q_4^{\mathsf T}=
 \begin{pmatrix}
 -A(u-1)H\\-B\bigl(AH(u-1)-1\bigr)\\
 -AuB(u-1)H\\-AuvB(u-1)H\\-cB(u-1)H
 \end{pmatrix},\qquad
 q_5^{\mathsf T}=
 \begin{pmatrix}
 -Aw(v-1)\\-AwB(v-1)\\-B\bigl(Auw(v-1)-1\bigr)\\
 -B(v-1)(c-1)\\-cwB(v-1)
 \end{pmatrix}.
 \end{gathered}                                                \tag{F.31.15}
\]
With rows in the exact order shown, the
four determinant identities are
\[
 \begin{aligned}
 \det(q_6,q_8,q_9,q_{10},q_1)
   &=cu^2vw(r-1)B^4(c-1)(rtv+uv-u),\\
 \det(q_6,q_8,q_9,q_{10},q_3)
   &=c^2u^2(r-1)B^4(w-1),\\
 \det(q_6,q_8,q_9,q_{10},q_4)
   &=rtc^2u^2(r-1)B^4(u-1),\\
 \det(q_6,q_8,q_9,q_{10},q_5)
   &=-cuvw(r-1)B^4(v-1)F,\\
 F&=(c-1)(rt+u)-cu/v.
 \end{aligned}                                                \tag{F.31.16}
\]
The certificate multiplies the original
word matrices, checks every entry of
(F.31.11), compares all eight rows with
(F.31.12)--(F.31.15), and expands each
determinant by its signed permutations.
The resulting identities hold in $A$,
before reduction in any characteristic.

\emph{The exact common fixed field and the generic comparisons.}
Let $\sigma(a)=a^{q_1}$ on $k_0$.
This automorphism has order exactly $m$:
its $m$th power is the identity, and for
$0<a<m$ the polynomial $X^{q_1^a}-X$
has degree smaller than $|k_0|$ and cannot
vanish on all of $k_0$.
Write $r_j=\sigma^j(r)$ and $t_j=\sigma^j(t)$.
The coefficient map $d\mapsto d^{q_1^j}$
is additive and semilinear, preserves the
actual cocycle rule and principal line,
commutes with each source precomposition,
and has inverse $d\mapsto d^{q_1^{m-j}}$.
It transports each exact basis map above
to $Q_j=H^1_{\rm cont}(V,k_{0,\chi_j})$
and its core $W_j$.
It also proves $k_0=k_1(r_j,t_j)$.

Take distinct $j,l$ and source-invariant
subspaces of $W_j,W_l$ with an intertwining
isomorphism. Source $0$ first forces
$r_j=r_l$, since it acts there by
$r_jI$ and $r_lI$.
Choose the unique $a\in\{1,\ldots,m-1\}$
congruent to $l-j$ modulo $m$, and define
the actual subfield
\[
       L_0=\{x\in k_0:\sigma^a(x)=x\}.
                                                               \tag{F.31.17}
\]
It contains $k_1$ and $r_j$.
The element $t_j$ lies outside $L_0$:
otherwise $\sigma^a$ would fix both
$r_j,t_j$ and $k_1$, hence their whole
field $k_0$, contradicting its proved
nonidentity. In particular
$t_j\ne t_l$ and $t_j-1\ne0$.

All the base roots in (F.31.10), evaluated
at the $j$th or $l$th character, coincide
and lie in $L_0$.
Each moving root is $\eta/t_j$ or
$\eta/t_l$ for the same nonzero
$\eta\in k_1$.
Thus the two moving roots are distinct,
both lie outside $L_0$, and every factor
of $\tau_{l',j}$ and $\tau_{l',l}$ is
nonzero, for every source $l'$ in that table.
This proves the exact coefficient-field
inclusions used in the following cancellation.

Identify the two subspaces through the
specified intertwiner. The common
source-$l'$ endomorphism satisfies the
two annihilators
$(Z-b_{l',j})Q_{l'}(Z)$ and
$(Z-b_{l',l})Q_{l'}(Z)$.
Subtracting gives
$(b_{l',l}-b_{l',j})Q_{l'}(Z)=0$ on it.
The nonzero scalar just proved shows
that $Q_{l'}(G_{l',j})$ kills the first
subspace. By (F.31.11) and
$\tau_{l',j}\ne0$, the row $q_{l',j}$
kills that subspace as well.
This argument permits all equalities
among the base roots, including every
characteristic-dependent coincidence.

The first four rows in (F.31.16) therefore
vanish on the subspace, as do all four
possible last rows.
At least one of those four square row
matrices is invertible. Indeed, for
$c\ne1$ use the first determinant:
its last factor cannot vanish, because
$r_jt_jv+uv-u=0$ would put
$t_j=(u-uv)/(r_jv)$ in $L_0$.
Its other factors are nonzero by
$r_j\ne1$, $t_j\notin L_0$, and $c\ne1$.
For $c=1$, choose the second determinant
when $w\ne1$, the third when $w=1$
and $u\ne1$, and the fourth when
$w=u=1$. In this last case $v\ne1$
because $C\ne\{1\}$, and its factor
$F$ is exactly $-u/v\ne0$.
All omitted prefactors in this discussion
are displayed Laurent monomial units.
The invertible row matrix forces the
subspace to be zero. This excludes any
nonzero common core subspace for distinct
Frobenius indices on the entire $r\ne1$ chart.

For a nonzero source-invariant subspace
of $Q_j$ with no proper nonzero invariant
subspace, its intersection with $W_j$
either equals the whole subspace or is
zero. In the latter case $\pi_j$ is an
injective intertwiner onto the boundary
character $\mu_j$.
A core subspace and a boundary subspace
cannot be isomorphic: their source-$0$
scalars are $r_j\ne1$ and $1$.
Two boundary characters at distinct
indices cannot coincide. Source $7$
would give $t_j=t_l$, and source $2$
would then give $r_j=r_l$. Since
$k_0=k_1(r_j,t_j)$, this would make
$\sigma^{l-j}$ the identity.
The inverse of an intertwining isomorphism
covers the reversed mixed comparison.
These are all three unordered types on
this chart, obtained through the exact
boundary quotient rather than an assumed
invariant complementary line.

\emph{The retained chart for $r=1$.}
Now $k_0=k_1(t)$ by (F.31.3).
The proper extension implies $t\notin k_1$,
so $t-1\ne0$.
Use the actual third principal coordinate:
\[
 \begin{aligned}
 P_3(z)&=(z_i-\delta_i(t-1)^{-1}z_3)_{i\in I},
       &I&=(1,2,4,5,6,7),\\
 S_3(x)&=(x_1,x_2,0,x_3,x_4,x_5,x_6)^{\mathsf T},\\
 P_3S_3&=I_6,
       &S_3P_3(z)&=z-\delta(t-1)^{-1}z_3.
 \end{aligned}                                                \tag{F.31.18}
\]
Thus $Q_\chi$ has the coordinates of the
original generators in the order $I$.
Since $\chi(X_1)=1$, evaluation at $X_1$
also descends through the principal line.
Put $b=(1,1,st,stu,stuv,c)$ and define
the exact two-coordinate quotient
\[
 \begin{gathered}
 \pi(z)=(bz,z_1),\qquad
 \iota(x_1,x_2,x_3,x_4)=
 (0,-stx_1-stux_2-stuvx_3-cx_4,x_1,x_2,x_3,x_4),\\
 \sigma(a,h)=(h,a-h,0,0,0,0),\qquad
 \pi\sigma=I_2,\qquad
 z=\iota(z_3,z_4,z_5,z_6)+\sigma\pi(z),\\
 J=(e_3-ste_2,e_4-stue_2,e_5-stuve_2,
                         e_6-ce_2,e_2,e_1-e_2),\\
 J^{-1}(z)=(z_3,z_4,z_5,z_6,bz,z_1).
 \end{gathered}                                                \tag{F.31.19}
\]
These formulas prove both basis inverses
and identify the invariant core
$W_\chi=\ker\pi$ with $k_0^4$.
Its quotient coordinates are the actual
boundary evaluation followed by the
actual $X_1$ evaluation.
The original-word identities are
\[
 \begin{gathered}
 T_l=J^{-1}P_3A_lS_3J=
 \begin{pmatrix}G_l&*\;&*\\0&\mu_l&0\\0&0&\nu_l\end{pmatrix},\\
 (t-1)(A_lS_3J-S_3JT_l)
     =\delta\,\operatorname{row}_3(A_lS_3J),\\
 (\nu_0,\ldots,\nu_{10})=
 (1,u^{-2},t^{-1},v^{-2},w^{-2},c^2,1,t,1,1,1).
 \end{gathered}                                                \tag{F.31.20}
\]
Here $\mu$ is (F.31.5) with $r=1$.
For each word the first row of $A_l$
is exactly $(\nu_l,0,\ldots,0)$;
this proves the $X_1$ quotient before
passing to cohomology. The cleared
identity proves the full block via
$P_3$ and the nonzero scalar $t-1$.

\emph{Four moving rows and all generic type comparisons for $r=1$.}
Use the first four sources of (F.31.10),
namely $l=6,8,9,10$, with their same
base and moving roots and $r=1$.
Their core matrices are now $4$ by $4$.
For the rows below they satisfy precisely
\[
 Q_l(G_l)(G_l-b_lI_4)=0,\qquad
 q_lG_l=b_lq_l,\qquad q_lQ_l(G_l)=Q_l(b_l)q_l.
                                                               \tag{F.31.21}
\]
The actual rows are
\[
 \begin{gathered}
 q_6^{\mathsf T}=
 \begin{pmatrix}
 -(st+1)(c-1)\\stu(u-1)(c-1)\\
 stu(c-1)(uv-u-v)\\c(cu-cu/v-c-u+1)
 \end{pmatrix},\qquad
 q_8^{\mathsf T}=
 \begin{pmatrix}1-w\\-u(w-1)\\uv\\0\end{pmatrix},\\
 q_9^{\mathsf T}=
 \begin{pmatrix}-stu+st+1\\-stu(u-1)\\-stuv(u-1)\\-c(u-1)\end{pmatrix},\qquad
 q_{10}^{\mathsf T}=
 \begin{pmatrix}1-v\\u\\0\\0\end{pmatrix},\\
              \det(q_6,q_8,q_9,q_{10})=u^2c^2.
 \end{gathered}                                                \tag{F.31.22}
\]
This determinant is a Laurent unit in
every characteristic and at every allowed
value of $u,v,w,c$.

All $t_j$ are distinct and lie outside
$k_1$, since $k_0=k_1(t)$ and the
$k_1$-Frobenius has order $m$.
The moving roots $\eta/t_j$ therefore
lie outside $k_1$, and differ for
different indices; all base roots lie
in $k_1$. Subtracting the annihilators
as above, then using (F.31.21) and its
nonzero $Q_l(b_l)$, shows that a common
core subspace is killed by all four
rows of (F.31.22). Its unit determinant
excludes every such nonzero subspace.

The same rows exclude any core subspace
isomorphic to a boundary or $X_1$
character from any Frobenius factor.
Such a subspace is a line, and for
each of the four negative sources its
scalar $\lambda_l$ lies in $k_1$:
the boundary values are
$(u^{-2},v^{-2},w^{-2},c^2)$ and
the $X_1$ values are all $1$.
The annihilator on the line gives
$(\lambda_l-b_{l,j})Q_l(\lambda_l)=0$.
The first factor is nonzero, so
$Q_l(G_{l,j})$ kills the line.
Equation (F.31.21) and $Q_l(b_{l,j})\ne0$
make all four rows vanish on it,
contradicting (F.31.22).
Inverse intertwiners cover both orientations.

A simple invariant subspace of $Q_j$
either lies in $W_j$, or projects
injectively under (F.31.19) to a simple
subspace of its two-dimensional diagonal
quotient. A nonzero projection of the
latter to either one-dimensional
coordinate is injective; its character
is therefore $\mu_j$ or $\nu_j$.
This argument uses only the displayed
coordinate maps and does not assume
semisimplicity of a larger source module.
For distinct indices, $\mu_j=\mu_l$
is excluded by source $7$, and
$\nu_j=\nu_l$ is excluded by source $2$.
These core comparisons and the two
equal-type comparisons leave only
the cross-type equality
$\mu_j=\nu_l$.
Sources $1,3,4,5$ force exactly
\[
                  u^2=v^2=w^2=c^2=1.
                                                               \tag{F.31.23}
\]
In characteristic two this would make
all four parameters equal to $1$ and
give $C=\{1\}$, already excluded.
Thus the remaining field has odd
characteristic, its four parameters
are signs not all $1$, $k_1=\mathbb F_p$,
and $t_j\notin\mathbb F_p$.

\emph{All fifteen sign strata: exact absence of a boundary eigenline.}
For each sign tuple write
$\epsilon=cuvw$.
The coefficient homomorphism is explicitly
\[
 \begin{aligned}
 \mathbb Z[s^{\pm1},t^{\pm1},u^{\pm1},v^{\pm1},w^{\pm1}]
    &\longrightarrow\mathbb Z[t^{\pm1}],\\
 s^at^bu^dv^ew^f
    &\longmapsto\epsilon^a u^dv^ew^f t^{b-a},
       \qquad s\longmapsto\epsilon/t,
 \end{aligned}                                                \tag{F.31.24}
\]
where on the right $u,v,w,c$ are the
specified integer signs.
Addition of exponent vectors multiplies
both their signs and their powers of
$t$, so this is an exact ring
homomorphism, including negative exponents.
It is applied to the complete actual
word matrices and all basis maps.

Let $Q$ have the four rows in (F.31.22)
after this map. Its entries are integers,
and $\det Q=u^2c^2=1$.
For $l=6,8,9,10$ let
$\tau_l=Q_l(b_l)$ and define an extended
six-coordinate row by the exact identity
\[
 \widetilde q_l=
       \tau_l^{-1}(q_l,0,0)Q_l(T_l),\qquad
 \widetilde q_l|_{k_0^4}=q_l.
                                                               \tag{F.31.25}
\]
The second equality follows from
(F.31.21), and its cleared integral
form is checked before any division.
Every base root here is a sign, while
each moving root is a sign times $t^{-1}$.
All factors of $\tau_l$ are nonzero
for $t\notin\mathbb F_p$.
The inverse in (F.31.25) is thus defined
in every remaining field, including
all odd characteristics.

On a boundary eigenline all four
negative-source scalars are $1$.
Consequently $Q_l(T_l)$, and hence
each $\widetilde q_l$, vanishes there.
The last two quotient scalars at
source $2$ are $t$ and $t^{-1}$.
They are distinct because $t\notin\mathbb F_p$.
A boundary eigenline projecting to
the boundary quotient therefore has
last coordinates $(a,0)$ with $a\ne0$;
the nonzero boundary coordinate also
follows from the already proved
exclusion of a boundary line in the core.
Scale by the unique $a^{-1}$ to obtain
last coordinates $(1,0)$.

If $E$ is the four-row matrix with
rows $\widetilde q_l$, its first four
columns are $Q$. The exact constraints
then force the first four coordinates
to be the unique vector
$x=-Q^{-1}E_{\mathrm{boundary}}$.
Both inverses of $Q$ are integral,
as $\det Q=1$.
The table gives this forced vector
and a nonzero coordinate of the
actual residual
$T_l(x,1,0)^{\mathsf T}-\mu_l(x,1,0)^{\mathsf T}$.
The source index $l$ is zero-based
as before; the coordinate $a$ in this
table is one-based in the six-coordinate
basis (F.31.19).
\[
\begin{array}{c|c|c|c|c}
 (u,v,w,c)&x&l&a&\text{residual coordinate}\\\hline
 (1,1,1,-1)&(0,0,0,0)&1&4&-2(t+1)/t\\
 (1,1,-1,1)&(0,0,-2,0)&3&3&2(t+1)/t\\
 (1,1,-1,-1)&(0,0,2,0)&1&3&-4\\
 (1,-1,1,1)&(0,0,0,0)&5&2&-2(t+1)/t\\
 (1,-1,1,-1)&(0,0,0,0)&1&2&4\\
 (1,-1,-1,1)&(0,0,-2,0)&3&2&4/t\\
 (1,-1,-1,-1)&(0,0,2,0)&1&2&4\\
 (-1,1,1,1)&(-2,0,0,0)&4&1&2(t+1)/t\\
 (-1,1,1,-1)&(-2,0,0,-4)&1&1&4\\
 (-1,1,-1,1)&(2,0,2,0)&3&1&4(2t-1)/t\\
 (-1,1,-1,-1)&(2,0,6,4)&1&1&4\\
 (-1,-1,1,1)&(-2,-4,0,0)&4&1&2(3t-1)/t\\
 (-1,-1,1,-1)&(-2,-4,0,-4)&1&1&-4\\
 (-1,-1,-1,1)&(2,4,2,0)&3&1&4(2t-1)/t\\
 (-1,-1,-1,-1)&(2,4,6,4)&1&1&-4
\end{array}.
                                                               \tag{F.31.26}
\]
All fifteen tuples other than
$(1,1,1,1)$ occur exactly once.
The constants $2,4$ and the parameter
$t$ are units in these fields.
The factor $t+1$ cannot vanish for
$t\notin\mathbb F_p$.
The factor $2t-1$ has nonzero leading
coefficient in every odd characteristic,
and its vanishing would put $t$ in
the prime field. For $3t-1$, characteristic
three gives the nonzero constant $-1$;
every other odd characteristic again
gives a prime-field root.
Thus every displayed residual is
nonzero in every remaining field.

The certificate checks both inverses of
each specialized $Q$, the cleared
identities (F.31.25), the four equations
$E(x,1,0)^{\mathsf T}=0$, and each
displayed residual directly in
$\mathbb Z[t^{\pm1}]$.
Uniqueness from $Q^{-1}$ makes the
table exhaustive for all possible
normalized boundary eigenvectors.
There is consequently no boundary
eigenline on any of the fifteen strata.
This excludes a pair of isomorphic
subspaces of types $\mu_j,\nu_l$;
the reverse orientation is excluded
by the inverse intertwiner.
Together with the earlier comparisons,
this accounts for all six unordered
type pairs $W/W$, $W/\mu$, $W/\nu$,
$\mu/\mu$, $\nu/\nu$ and $\mu/\nu$.

\emph{Summed restriction through the actual invariant maps.}
All characters $\chi_j$ restrict to
the same character on $K$, because
$\chi(K)=C\subset k_1^\times$.
The actual restriction therefore defines
\[
 \Phi:\bigoplus_{j=0}^{m-1}Q_j
  \longrightarrow H^1_{\rm cont}(K,k_{0,\chi|K}),\qquad
       ([d_j])_j\longmapsto\sum_j[d_j|K].
                                                               \tag{F.31.27}
\]
It commutes with each of the eleven
source operators, since their actual
automorphisms preserve $K$ and all
the characters. Every individual
restriction is injective by F.24,
as $\chi_j(R)=C\ne\{1\}$.

Suppose that $\ker\Phi$ were nonzero,
and choose inside it a nonzero invariant
subspace $T$ of smallest dimension.
Finite dimension makes this choice
available. Each nonzero coordinate
projection $p_j|T$ is injective:
its kernel is an invariant subspace
of $T$, and cannot be all of $T$.
Its image has no proper nonzero
invariant subspace, by pulling such
a subspace back through this inverse.
Two nonzero coordinate projections
give the explicit mutually inverse
intertwiners
\[
 \begin{gathered}
 p_l\circ(p_j|T)^{-1}:p_j(T)\longrightarrow p_l(T),\\
 p_j\circ(p_l|T)^{-1}:p_l(T)\longrightarrow p_j(T).
 \end{gathered}                                                \tag{F.31.28}
\]
The complete comparisons on the
$r\ne1$ and $r=1$ charts exclude
such isomorphisms for $j\ne l$.
Hence exactly one projection is
nonzero. The subspace $T$ then lies
in one factor, where $\Phi$ is the
injective individual restriction,
a contradiction. This proves the
summed $H^1$ injection.

Zero summed restriction of cocycles
now forces every cocycle to be
principal. Since the common restricted
character has nontrivial image $C$,
the exact kernel is
\[
 \left\{\bigl(b_j(\chi_j-1)\bigr)_{j=0}^{m-1}:
                      \sum_{j=0}^{m-1}b_j=0\right\}.
                                                               \tag{F.31.29}
\]
Its converse inclusion follows directly
by restriction. Each principal vector
is nonzero, so this space has dimension
$m-1$. The cocycle domain has dimension
$7m$ by its actual generator coordinates,
and the claimed rank is $6m+1$.

For one cocycle $d$, apply this result
to the actual coefficient tuple
$(d^{q_1^j})_{j=0}^{m-1}$.
It gives the exact trace-kernel identity
\[
 \operatorname{Tr}_{k_0/k_1}(d|K)=0
 \quad\Longleftrightarrow\quad
 d=b(\chi-1),\qquad
                    \operatorname{Tr}_{k_0/k_1}(b)=0.
                                                               \tag{F.31.30}
\]
Indeed the forward implication makes
$d$ principal. Its Frobenius powers
then have the unique parameters
$b^{q_1^j}$, whose sum is the trace;
(F.31.29) gives the stated vanishing.
The reverse implication uses
$\chi(K)\subset k_1$ and the trace's
$k_1$-linearity.

The explicit conjugation, translation
projection and trace-pairing coefficient
maps in F.26 send every proper finite
affine marginal image to a nonprincipal
element of this trace kernel over the
actual field $k_0$.
Those maps retain the scalar character
and both scalar fields, so they retain
(F.31.1). Equation (F.31.30) excludes
the resulting nonprincipal element.
This proves (F.31.2) for the negative
marginal, for every original finite
coefficient field including fields
larger than $k_0$.

\emph{Integral certificate, continuity and sign transport.}
The independently structured verifier
imports none of the four new producers
and no symbolic algebra package.
It reconstructs all eleven literal
decorated Artin words on each of the
two coefficient charts using the
previously reviewed integer Laurent
word arithmetic. It checks the full
seven-coordinate matrices and the
quotient blocks by the cleared
intertwining identities (F.31.9) and
(F.31.20), independently of the
producer's Laurent-division algorithm.
Its $1603$ exact checks include every
displayed basis inverse, all twelve
moving-root annihilators and direct
left-eigenrow identities, the five
row determinants, and all fifteen
coefficient maps, forced lifts and
residual obstructions.
The nonmonomial principal pivots are
exactly $r-1$ and $t-1$ on their
respective charts. Their nonvanishing,
and every further factor cancelled
in the proof, were established in
the actual finite fields before use.
The evidence manifest binds the
source, programs, original words,
coefficient inputs, outputs and
completed dependencies.

Every cocycle constructed from generator
values is the unique continuous extension
of a homomorphism to a finite affine
group. The source automorphisms are
continuous by F.23, and all maps between
the finite coefficient fields are
continuous. Their cocycle compositions
and the restrictions to the actual
closed subgroup $K$ are therefore
continuous throughout the argument.
For the positive marginal, precompose
with the actual based isomorphism
$\mathcal B_+:V_-\to V_+$ of F.23.13.
Its inverse is supplied there, and
it sends $X_i$ precisely to $Y_i^+$
and carries $\mathcal K_-,R_-$ onto
their positive counterparts.
It therefore preserves the seven
specified character values, their
product, both scalar fields and
all maps defining the summed restriction.
The proved negative image equality
and kernel statements give the
positive statements through these
explicit mutually inverse maps.
This proves the claim for both signs.
\end{proof}
